
Our work bridges measurement theory from psychometrics with reinforcement learning reward design for code generation. Existing research has advanced execution-based RL, proxy-based alignment, and evaluation metrics independently, but no prior work systematically validates proxy construct validity before optimization.

\subsubsection{Execution-Based RL for Code Generation}

CodeRL \citep{Le2022CodeRL} introduced actor-critic training with unit test pass/fail outcomes as reward signals, achieving significant improvements over supervised fine-tuning on HumanEval (pass@1: 28.8\% $\rightarrow$ 36.8\%). CURE \citep{Tian2023} extended this through co-evolutionary training of code generators and test generators, achieving 39.6\% pass@1. DRIVE-RLVR \citep{Liu2024} introduced data curation practices for reinforcement learning with verifiable rewards in competitive programming.

These works demonstrate that execution feedback is a reliable reward signal but do not explore whether auxiliary proxies (structural similarity, efficiency, style) can augment execution rewards. None report measurement reliability metrics (CV, test-retest correlation, cross-platform stability) for auxiliary reward signals. This gap motivates our contribution: existing work adopts execution and auxiliary rewards without pre-validation.

\subsubsection{Proxy-Based Alignment for Code Generation}

CodeUltraFeedback pioneered reinforcement learning from human feedback for code generation using LLM-as-a-judge to generate preference rankings, achieving performance gains on CODAL-Bench (HumanEval: 72.3\% $\rightarrow$ 78.1\%). SEAlign extended alignment to multi-step repository-level tasks through Monte Carlo Tree Search, achieving state-of-the-art on SWE-Bench (12.6\% resolution rate). SelfCodeAlign demonstrated alignment without human annotations through self-generated instruction-response pairs.

However, these works do not test conditional independence: do quality assessments from LLM judges explain variance in behavioral outcomes after controlling for execution pass rate? If preference signals are conditionally redundant---correlating with outcomes only because they proxy for execution correctness---then alignment optimizes a noisier version of the execution signal rather than an orthogonal quality dimension. Our Stage 2 validation (hierarchical regression testing $\Delta R^2 \geq 0.03)$ addresses this gap.

\subsubsection{Code Generation Evaluation Metrics}

Chen et al. (2021) introduced CodeBLEU alongside the HumanEval benchmark, validating the metric's correlation with human judgments of code quality ($\rho$=0.52). ExeDS introduced behavioral correctness testing via Jupyter notebook execution, demonstrating that surface-form metrics poorly predict execution correctness (Spearman $\rho$=0.23). RepoBench established repository-level code completion evaluation. AutoGEEval++ extended evaluation to domain-specific code generation with multi-dimensional metrics including runtime efficiency.

These works validate metrics via correlation with human judgment (validity testing) or demonstrate that metrics capture distinct aspects of quality. Our contribution is reliability testing: are metrics stable (CV), discriminative (Cohen's d), and generalizable (Spearman $\rho$)? These are prerequisite questions for optimization---a noisy metric may correlate with quality yet be unsuitable for RL reward functions.

\subsubsection{Multi-Objective RL and Structured Rewards}

Lei Chen et al. (2025) introduced multi-granularity structured RL for chart-to-code generation using dual reward signals (textual similarity and visual similarity), demonstrating that MSRL breaks the supervised fine-tuning plateau (HumanEval: 68.2\% SFT $\rightarrow$ 74.5\% MSRL). This result supports the hypothesis that multi-dimensional optimization can yield improvements, but the work does not test measurement reliability of the visual reward signal or conditional independence of reward components.

Our work differs by introducing systematic construct validation before optimization. If proxies fail validation, we identify issues before multi-objective training. This prevents scenarios where RL algorithms expend computational resources balancing noisy or redundant reward components.

\subsubsection{Summary and Positioning}

Existing execution-based RL demonstrates that test-driven feedback is reliable but does not explore auxiliary dimensions. Existing proxy-based alignment adopts auxiliary metrics without pre-validation. Existing evaluation metrics research validates correlation with quality but not measurement stability. Existing multi-objective RL shows multi-granularity rewards improve performance but does not test construct validity.

We fill this gap by introducing a four-stage validation pipeline with Stage 1 empirically validated. Our methodological contribution is the framework itself---reusable across code quality dimensions and applicable beyond code generation to any proxy-based optimization domain. Our empirical contribution demonstrates the framework's functionality: CodeBLEU passes Stage 1 with exceptional margins, while runtime efficiency using wall-clock measurements marginally exceeds the CV threshold, suggesting that hardware instrumentation may be required. This compositional finding validates the four-stage pipeline's design: test each proxy independently before multi-objective optimization.
